\section{Forecasting Task and Benchmark}
\label{sec:benchmark}

We study LLM forecasting as a prospective systems task: a model must issue a timestamped judgment
about an event that has not yet occurred, express its uncertainty in a structured form, and do so
under a precisely defined information condition. We instantiate this task at the 2026 FIFA World Cup.
The tournament combines a fixed universe of 104 matches and objectively resolvable outcomes with the
features that make real-world forecasting difficult: information changes over time, evidence can be
incomplete or contradictory, and low-scoring matches retain substantial irreducible uncertainty.
Football is therefore a controlled test bed for forecasting mechanisms rather than a claim that the
resulting rankings transfer unchanged to every decision domain.

\subsection{Platform and Forecasting Task}
\label{sec:platform}

\arena is both a prospective forecast-collection system and a public interface to the resulting
benchmark archive. Before each match, the platform records forecasts from contemporary LLMs and makes
them inspectable before the outcome is known. Once an official result becomes available, the same
stored forecast is scored and incorporated into the public leaderboard and the research dataset.
Match pages show the predicted scoreline, home/draw/away probabilities, advancement probabilities for
knockout matches, the forecast condition, timing, confidence, and a short rationale. The analytics
view supports comparisons by model, metric, information access, prompt strategy, horizon, tournament
stage, and date; separate views display the 15 tournament questions and the knockout bracket. Website
visits and public interaction are not experimental units and do not affect the forecasts.

The universal match-level target is the official result after regulation time plus stoppage time.
For match $i$, each model returns a probability vector
$\mathbf{p}_i=(p_{i,\mathrm{H}},p_{i,\mathrm{D}},p_{i,\mathrm{A}})$ for a home win, draw, and away win,
respectively. Extra time and penalty shootouts are excluded from this target. Knockout fixtures
therefore receive an additional two-class forecast for which team advances. Each response also
contains expected home and away goals, one most likely 90-minute scoreline, full-match fields,
self-reported confidence in $[0,1]$, and a short rationale. Before the tournament, the same model
configurations answered 15 longer-horizon questions: the 12 group winners, the set of four
semifinalists, the world champion, and the team of the tournament's top scorer.

Figure~\ref{fig:architecture} summarizes the prospective data flow. A fixture registry supplies only
match-identifying information: competition, stage, UTC kickoff time, listed home and away teams,
venue when known, and knockout status. A scheduler expands each event into the prespecified
experimental cells. The prompt builder combines the shared fixture block and output schema with the
assigned access and elicitation instructions. After model execution, responses pass through schema,
range, probability-sum, and scoreline validation. Raw prompts, raw responses, tool traces, and all
validation outcomes remain stored alongside the parsed forecast. Official outcomes cross the
pipeline only after the forecast has been committed. The website and the reproducible analysis read
from this common archive; neither feeds information back into subsequent forecasts.

\begin{figure*}[t]
  \centering
  \begin{tikzpicture}[
    x=1cm,y=1cm,>=Latex,
    flow/.style={-Latex,thick,draw=draftblue},
    postflow/.style={-Latex,thick,draw=todoorange},
    prebox/.style={draw=draftblue,rounded corners=2pt,fill=draftlight,
      align=center,font=\scriptsize,minimum height=1.45cm,text width=2.40cm,inner sep=3pt},
    postbox/.style={draw=todoorange,rounded corners=2pt,fill=todolight,
      align=center,font=\scriptsize,minimum height=1.40cm,text width=2.55cm,inner sep=3pt}
  ]
    \node[prebox] (registry) at (0,2.25) {\textbf{Event registry}\\fixtures + 15 questions\\teams, stage, time, venue};
    \node[prebox] (scheduler) at (3.05,2.25) {\textbf{Factorial scheduler}\\8 models $\times$ 3 horizons\\$\times$ 2 access $\times$ 2 prompts};
    \node[prebox] (prompt) at (6.10,2.25) {\textbf{Prompt builder}\\shared target definitions\\condition blocks + schema};
    \node[prebox] (runtime) at (9.15,2.25) {\textbf{LLM execution}\\tools disabled/enabled\\temperature 0; $n=1$};
    \node[prebox] (validation) at (12.20,2.25) {\textbf{Validation}\\JSON, fields, ranges, sums\\normalization or repair};
    \node[prebox] (archive) at (15.30,2.25) {\textbf{Timestamped archive}\\raw + validated forecast\\confidence, rationale, traces};

    \draw[flow] (registry) -- (scheduler);
    \draw[flow] (scheduler) -- (prompt);
    \draw[flow] (prompt) -- (runtime);
    \draw[flow] (runtime) -- (validation);
    \draw[flow] (validation) -- (archive);

    \node[prebox] (public) at (15.30,0.05) {\textbf{Public platform}\\pre-match predictions\\analytics, questions, bracket};
    \draw[flow] (archive) -- (public);

    \draw[dashed,thick,draw=todoorange] (-1.35,-1.02) -- (16.75,-1.02);
    \node[fill=white,inner sep=2pt,font=\scriptsize\bfseries,text=todoorange]
      at (7.7,-1.02) {Outcome barrier: forecasts are stored before resolution};

    \node[postbox] (outcome) at (0,-2.70) {\textbf{Official resolution}\\90-minute score, advancer\\tournament answers};
    \node[postbox] (freeze) at (4.10,-2.70) {\textbf{Freeze and reconcile}\\consistent database backup\\integrity check + hashes};
    \node[postbox] (score) at (8.20,-2.70) {\textbf{Evaluation}\\proper scores + accuracy\\calibration + diagnostics};
    \node[postbox] (release) at (12.30,-2.70) {\textbf{Research release}\\typed tables + manifest\\figures and audit logs};

    \draw[postflow] (outcome) -- (freeze);
    \draw[postflow] (freeze) -- (score);
    \draw[postflow] (score) -- (release);
    \draw[postflow] (archive.south) |- (freeze.north);
    \draw[postflow] (release.east) -| (public.south);
  \end{tikzpicture}
  \caption{Architecture and prospective data flow of \arena. The upper lane records forecasts and
  their provenance before resolution; the lower lane begins only after official outcomes become
  available. The public platform and research pipeline remain linked to the same forecast records.}
  \label{fig:architecture}
\end{figure*}

The public interface is shown separately in Figure~\ref{fig:website}. This separation distinguishes
the benchmark architecture from its presentation layer: Figure~\ref{fig:architecture} explains how
forecasts are created and protected from outcome leakage, whereas Figure~\ref{fig:website} shows how
individual predictions and aggregate analyses can be inspected.

\begin{figure*}[t]
  \centering
  \begin{subfigure}[t]{0.485\textwidth}
    \centering
    \includegraphics[width=\linewidth]{figures/screenshot-card.png}
    \caption{Match forecasts, probabilities, conditions, confidence, and rationales.}
  \end{subfigure}\hfill
  \begin{subfigure}[t]{0.485\textwidth}
    \centering
    \includegraphics[width=\linewidth]{figures/screenshot-analytics.png}
    \caption{Filterable leaderboard and performance trajectories.}
  \end{subfigure}
  \caption{Public views over the benchmark archive. Forecast details remain visible at the level of
  individual matches, while the analytics interface compares models and experimental conditions.}
  \label{fig:website}
\end{figure*}

Although the present deployment focuses on football, the platform is deliberately modular. Its event
registry, scheduler, prompt components, outcome resolver, and scoring layer are separated, so that a
new target can be added by providing a domain-specific event adapter, forecast schema, resolution
rule, and evaluation module. The same infrastructure can therefore support other sports and live
competitions, as well as non-sporting domains with timestamped and objectively resolvable future
events.

\subsection{Data Collection and Experimental Dimensions}
\label{sec:data-collection}

The benchmark deliberately varies more than model identity. Eight contemporary proprietary and
open-weight LLMs---including GPT-5.5, Gemini~3.1~Pro, and Claude~Opus~4.8---are evaluated under a
crossed design of information access, elicitation strategy, and forecast horizon. Exact API names,
canonical versions, providers, model families, tool support, and primary-panel eligibility are
reported in one appendix table. All complete prompts and JSON schemas are likewise moved to the
appendix. Calls use one deterministic sample per experimental cell with temperature zero.

\begin{table*}[t]
  \caption{Dimensions captured by the benchmark. The design varies deployment conditions while
  retaining a common fixture block and response schema.}
  \label{tab:benchmark-dimensions}
  \centering
  \small
  \begin{tabularx}{\textwidth}{>{\raggedright\arraybackslash}p{0.18\textwidth}
    >{\raggedright\arraybackslash}p{0.30\textwidth}X}
    \toprule
    \textbf{Dimension} & \textbf{Levels or fields} & \textbf{Analytical purpose} \\
    \midrule
    Model & Eight LLMs from seven provider families &
    Compares contemporary systems while separating model identity from deployment choices. \\

    Information access & Closed book; open book with web search enabled &
    Tests the value of current public information. Assigned access, observed search, number of tool
    calls, and trace availability are stored separately so that search behavior can be analyzed. \\

    Prompt strategy & Score first; probabilities first &
    Reverses elicitation order while keeping target definitions and output fields fixed, enabling a
    controlled test of whether prompt order changes quality or forecast structure. \\

    Forecast horizon & Stage opening; T--24h; T--2h &
    Repeats forecasts as information develops. Scheduled and actual lead time are both retained;
    T--24h is the primary operational snapshot. \\

    Structured forecast & H/D/A probabilities, expected goals, modal scoreline, advancement &
    Separates probabilistic quality, categorical decisions, score prediction, and knockout progress. \\

    Confidence & Overall self-reported confidence in $[0,1]$ &
    Tests whether a model's self-assessment ranks relatively easier and harder forecasts; confidence
    is not assumed to be a calibrated probability of correctness. \\

    Observable reasoning & Short rationale plus provider search/tool trace &
    Supports qualitative analysis of which evidence types are mentioned and whether assigned search
    is observably executed, without claiming access to private chain of thought. \\

    Output reliability & Raw response, validation, normalization, repair, and error status &
    Measures whether a forecasting system produces a valid, coherent, and scoreable response rather
    than conditioning all comparisons on successful calls. \\
    \bottomrule
  \end{tabularx}
\end{table*}

Closed-book calls receive the fixture block but no external tools and are instructed not to browse.
Open-book calls receive the identical fixture block, are given web-search access, and are instructed
to check current information before forecasting. Tool availability is not equated with tool use:
observable search calls and traces are logged, and the primary access comparison follows an
intent-to-treat definition based on assigned condition. The score-first prompt requests the modal
scoreline before probabilities and expected goals; probabilities-first reverses this order while
requiring the same fields. Forecasts are repeated at stage opening, approximately 24 hours before
kickoff, and approximately two hours before kickoff.

This design produced 9,984 match-level forecast attempts
($104\times8\times3\times2\times2$), of which 9,316 are valid for scoring. Seven models provide a
complete 1,248-record panel each; the eighth model contributes 580 valid forecasts and 668 retained
API errors. The 15 tournament questions add 480 forecasts, yielding 10,464 scheduled primary
forecasting prompts overall, before any schema-repair or transport-level retry calls. Every successful
provider response retains its billed cost, allowing the final project expenditure to be reconstructed
from the frozen archive; the invoice-reconciled total is
USD~\TODO{insert final all-in amount}. As a concrete reference, open-book T--24h execution costs on
average USD~0.110 more per forecast than the matched closed-book condition. We report expenditure
once to communicate the scale of collection; cost and latency are logged for reproducibility but are
not central forecast-quality outcomes.

The reasoning analysis combines objective tool metadata with generated rationale text. For each
open-book attempt, we distinguish observed search, no observed search, and unknown observability.
Non-empty rationales are coded independently by a released keyword lexicon and a fixed-prompt LLM
annotator. All disagreements and a fixed-seed sample of agreements are exported without match,
model, outcome, or performance identifiers for blinded human adjudication. Final category prevalence
therefore describes observable mentions of evidence---such as form, injuries, lineups, rankings,
market information, or tournament context---rather than hidden reasoning, factual correctness, or
proof that a mentioned source was actually consulted.

\subsection{Evaluation and Benchmark Strengths}
\label{sec:benchmark-strengths}

The primary outcome measure is the unscaled three-class Brier score,
$\mathrm{BS}_i=\sum_{c\in\{\mathrm{H,D,A}\}}(p_{i,c}-y_{i,c})^2$, where zero is optimal and a
uniform forecast scores $2/3$. Log loss complements the Brier score by penalizing confident errors
more strongly. We additionally report probability-modal H/D/A accuracy, exact-score accuracy,
goal-difference and tendency accuracy, absolute goal errors, and Kicktipp points (five for an exact
score, two for the correct goal difference, one for the correct tendency, and zero otherwise).
Knockout predictions receive advancement Brier score, log loss, and accuracy. Calibration is examined
through outcome-specific reliability diagrams; confidence is analyzed by within-configuration ranks.
Further diagnostics cover agreement between the score-implied outcome and probability argmax,
probability-sum constraints, pairwise probability correlation, Jensen--Shannon divergence,
same-cell ensembles, output validity, observed search, and qualitative rationale categories
\citep{brier_1950,gneiting_raftery_2007}.

\begin{table*}[t]
  \caption{Metric families used to evaluate the forecasts. No single measure is treated as a complete
  definition of forecasting quality.}
  \label{tab:evaluation-metrics}
  \centering
  \scriptsize
  \begin{tabularx}{\textwidth}{>{\raggedright\arraybackslash}p{0.19\textwidth}
    >{\raggedright\arraybackslash}p{0.32\textwidth}X}
    \toprule
    \textbf{Family} & \textbf{Measures} & \textbf{What they evaluate} \\
    \midrule
    Probability quality & Brier score; log loss; advancement scores &
    Quality of the full probability vector, including the severity of confident errors. \\
    Decision and score quality & Modal accuracy; exact score; goal difference; tendency; goal errors; Kicktipp points &
    Correct categorical decisions and football-specific scoreline performance. \\
    Calibration and confidence & Outcome reliability; within-cell confidence association &
    Agreement between probabilities and empirical frequencies, and whether confidence orders
    relative forecast difficulty. \\
    Coherence and diversity & Modal agreement; probability-sum audits; correlations; JS divergence; ensemble gain &
    Compatibility of structured fields and the extent to which models provide distinct signals. \\
    Reliability and tool use & Validity; repair/error rate; search-observed rate; rationale coverage and categories &
    Operational completeness, compliance with assigned search, and observable evidence content. \\
    Tournament questions & Choice Brier/log loss; pick accuracy; exact-set and team recovery &
    Longer-horizon single-choice and fixed-set forecasting beyond individual matches. \\
    \bottomrule
  \end{tabularx}
\end{table*}

The benchmark has four design strengths. First, it is genuinely prospective: every prediction is
timestamped before the corresponding outcome exists, reducing both training-data contamination and
post-outcome reconstruction. Second, the factorial design separates model choice from information
access and prompt strategy while holding the task and output schema fixed. Third, repeated horizons
permit within-match analysis of whether forecasts change as resolution approaches. Fourth, the
archive connects public inspectability with scientific provenance by retaining failed calls, raw and
validated responses, confidence, rationales, tool behavior, and timing. The platform can therefore
analyze not only which system forecasts best, but also how a forecast was elicited, whether search was
used, how uncertainty was expressed, and where a structured response failed.

\subsection{Research Questions and Statistical Analysis}
\label{sec:research-questions}

The multidimensional design supports the following analyses:

\begin{enumerate}[leftmargin=1.6em,itemsep=2pt,topsep=2pt]
  \item[\textbf{RQ1}] Does open-book information access improve forecast quality, and do gains vary
  by model, stage, or forecast horizon?
  \item[\textbf{RQ2}] Does score-first versus probabilities-first elicitation change forecast quality,
  calibration, scoreline structure, or cross-field coherence?
  \item[\textbf{RQ3}] How well calibrated are the elicited H/D/A probabilities, and does
  self-reported confidence contain information about subsequent forecast accuracy?
  \item[\textbf{RQ4}] How diverse are model forecasts within identical conditions, does open-book
  access produce convergence, and how much does same-cell ensembling improve robustness?
  \item[\textbf{RQ5}] When search is enabled, how often is it observably used, and how do the evidence
  categories mentioned in rationales differ between closed- and open-book forecasts and across models?
  \item[\textbf{RQ6}] How well do models answer longer-horizon tournament questions, and are their
  single-choice and fixed-set probabilities structurally coherent?
\end{enumerate}

T--24h and the 90-minute H/D/A target define the primary analysis. Primary model comparisons use the
seven-model complete panel; invalid forecasts are not imputed, and the partial eighth model is
reported through reliability statistics and labeled common-match sensitivity analyses. Because the
same match is forecast repeatedly, calls are not treated as independent observations. Prespecified
contrasts first aggregate nuisance dimensions within match and then form one paired difference per
match. Uncertainty uses 10,000 studentized bootstrap replicates that resample whole matches within
group-stage and knockout strata. Two-sided randomization inference uses 10,000 sign flips of the same
paired differences, and Holm correction is applied within declared comparison families. Reported
results include the estimand, 95\% interval, raw and adjusted $p$-values where applicable, and both
match and forecast counts. Tournament questions use the question---not each of its 32 responses---as
the unit of generalization. Confidence, rationale, and embedding analyses remain explicitly
exploratory and are not interpreted as causal evidence about hidden reasoning.
